% Add "answers" option to print the answers
\documentclass[addpoints, answers]{exam}

\usepackage[utf8]{inputenc}
\usepackage{calc}
\usepackage[dvipsnames]{xcolor}
\usepackage{ifthen}
\usepackage{tcolorbox}

% \makebox[\textwidth]{Name:\enspace\hrulefill}
%  
% \vspace{5mm}
%\makebox[\textwidth]{Instructor’s name:\enspace\hrulefill}

\definecolor{MyLightGray}{gray}{0.60}

\newcommand{\truefalse}
{		\begin{checkboxes}
			 \choice True
			 \choice False. Explanation: 
		\end{checkboxes}
		\vspace{10pt}
}


% Check if 'answers' is an option
\newcommand{\ifanswers}[2]{%
  \ifthenelse{\boolean{exam@answers}}{#1}{#2}%
}



\newcommand{\answerspace}[1]
{
	\ifprintanswers
	\else
		\setlength{\mylength}{10pt}%
		% \vspace{#1\mylength}
		% Make the grey box
		\par\nobreak\noindent
		\makebox[\linewidth][l]{%
			% set inner padding to zero, so only the frame shows
			{\setlength{\fboxsep}{0pt}%
			% draw a color frame, white (or transparent) inside
			\fcolorbox{MyLightGray}{white}{%
			% the inner box has zero width content but fixed height
			\parbox[t][#1\mylength][c]{%
				% subtract twice the frame‐thickness from the full width
				\dimexpr\linewidth-2\fboxrule\relax
			}{\strut}%
			}%
			}%
		}%
		\par
		\vspace{3\mylength}

  	\fi
}
\newlength{\mylength}

\newcommand{\code}[1]{{\texttt{#1}}}

\newcommand{\answer}[1]{\begin{solution}{#1}\end{solution}}




%\usepackage{listings}
\usepackage{color}
\usepackage{lastpage}
\usepackage{verbatim}
\usepackage{fancyvrb}
\usepackage{graphicx}
\usepackage{float}
\usepackage{minted}

\newcommand{\includecode}[2][c]{\lstinputlisting[caption=, escapechar=, language=#1]{#2}<!---->}
\RecustomVerbatimCommand{\VerbatimInput}{VerbatimInput}{}

\renewcommand{\solutiontitle}{\noindent\textbf{Solution sketch: }}

\definecolor{SolutionColor}{rgb}{0.8,0.9,1}

\cfoot{Page \thepage\ of \pageref{LastPage}}

\rhead{Initials: \rule{20mm}{0.4pt}}



%%%%%%%%%%%%%%%%%%%%%%%%%%%%%%%%%%%%%%%%
\begin{document}
\fvset{numbers=left,numbersep=3pt}
\thispagestyle{empty}

\begin{center}
	{\LARGE CS 408 Spring 2026 -- Midterm (A)\\}

	\vspace{20mm}
\fbox{\fbox{\parbox{5.5in}{\centering
	\begin{itemize}
          % \item Write your \textbf{name} in the first page and your \textbf{initials in every other page}.
        \item The exam is closed notes, closed book.

        \item The exam should be completed individually: any
          communication with your fellow students is strictly
          prohibited.

	\item Please print or write legibly.

        \item If you have a question, raise your hand and the
          course staff will come to you.

        \item Provide complete and clear but concise
          answers. Justify thoroughly statements where
          required.

        \item Answer the questions in the spaces provided. Your answer
          does not need to take up all the space provided; if you run
          out of room, use the empty pages at the back of the exam to
          complete your answer.

        \item Coding questions should be answered in Java, unless
          requested otherwise.

        \item Unless otherwise noted, CFG nodes represent basic blocks.
          % or low-level pseudo-code.
          % \item Unless otherwise noted, assume that questions refer to C code, AT\&T assembly syntax and the x86 architecture.

        \item Please double check that your copy of the exam
          is complete; you should have \pageref{LastPage}
          pages.

        \item You have 75 minutes to complete the exam.
        \end{itemize}
	}}}
\end{center}


\vspace{13mm}
\begin{center}
	\makebox[0.6\textwidth]{
		\combinedgradetable[h][questions]
	}
\end{center}
%\begin{center}\combinedgradetable[h][questions]\end{center}
\vspace{28mm}


\makebox[\textwidth]{\textbf{Name (please print)}: \rule{100mm}{0.8pt} }
%\end{center}

\vspace{10mm}
%\begin{center}
\makebox[\textwidth]{\textbf{PUID}: \rule{40mm}{0.8pt} \hspace{35mm}}
%\end{center}




\begin{questions}

  \newpage
  % \question State whether the following statements are true of false. \textbf{Explain the false statements.}
  %   \begin{parts}
  %   \part[2]
  %     \truefalse
  %     \answer{True.}
  %
  %   \end{parts}

\question[4] What are the key differences between black-box
  and white-box testing approaches?

  \answerspace{9}
  \answer{
    An error is an incorrect internal state of the program, such as a wrong value for a variable.
  }

\question[4] List some of the advantages and disadvantages of
  both black-box and white-box testing.

  \answerspace{9}
  \answer{
    An error is an incorrect internal state of the program, such as a wrong value for a variable.
  }


\question Consider the following control-flow graph:
	\begin{figure}[ht]
          % \includegraphics[width=7cm]{midterm-cfg1.pdf}
    %         public static void process(int n) {
    %     int i = 0;

    %     while (i < n) {

    %         // Meaningful but iteration-invariant condition
    %         if (n % 2 == 0) {
    %             System.out.println("Even n");
    %         } else {
    %             System.out.println("Odd n");
    %         }

    %         i++;
    %     }

    %     System.out.println("Done");
    % }

	\end{figure}

\vspace{0.2cm}

\begin{parts}
\part[3] If possible, provide 2 valid test paths through this program
  \textit{and} the test case inputs that would generate these
  paths. If this is not possible, explain why.

	\answerspace{7}
	\answer{
		Examples: [1, 3, 4, 5] and [1, 2, 4, 5].
		The test paths start from node 1 and end at node 5. There can be many possibilities of test paths in the given cfg, out of which, one path can be [1, 3, 4, 5] and another can be [1, 2 , 4, 5].
	}

	\part[3] If possible, list 3 simple paths that (a) are not prime paths and (b) start on node 3. If this is not possible, explain why.
	\answerspace{7}
	\answer{
		[2],
		[2, 4],
		[2, 4, 5],
		[2, 4, 1]
		\\
		Any three of the above four simple paths are acceptable answers.
	}

	\part[5] If possible, list four prime paths that start on node 4. If this is not possible, explain why.
	\answerspace{7}
	\answer{
		There are only 2 prime paths that start on 4 (which are [4, 1, 2, 4] and [4, 1, 3, 4]);  in total there are only 6 simple paths starting on 4, out of which only these two are not the proper subsets of other simple paths.
		}


	\newpage
      \part[5] Write an example of a syntactically valid Java function that generates
	this control-flow graph. If this is not possible,
	explain why.
	\answerspace{23}
	\answer{
		This is the closest possible solution-
		% \VerbatimInput{midterm-q7_d_answer.txt}
	}

\end{parts}

\newpage
\question Consider the following Java program:
  \begin{minted}{java}
public static void process(int x) {
  int y = Math.abs(x);
  while (y >= 0) {
    if (y < 0) {
       System.out.println("Negative y");
     }
     y--;
  }
  return z;
}
\end{minted}


\begin{parts}

\part[4] Is it possible to devise a set of test cases that achieve
  100\% \textbf{node coverage} for \mintinline{java}{process}?  %
  If your answer is yes, provide a minimal set of test cases that do
  so. %
  If your answer is no, (a) provide a simple change to
  \mintinline{java}{process} so that it becomes possible to satisfy
  100\% \textbf{node coverage} and (b) provide a simple and minimal
  test case set that satisfies \textbf{node coverage} of this
  function.  \answerspace{14}
  \answer{\\
        % rubric:
        % Case 1 - Yes:
        %       - explain: 2
        %       - test case: 1
        % Case 2 - No:
        %       - explain: 1
        %       - change: 1
        %       - test case: 1
        No. Because within the while loop, \texttt{x} will never be less than zero due to \texttt{unsigned int} type.
        Thus, the node, \texttt{line 7}, will never be covered. \\
        Example change: To satisfy node coverage with a simple change, changing \texttt{unsigned int x} to \texttt{int x}, and adding \texttt{x -= 1;} after line $5$ in the while loop.
        In this case, the test case, \texttt{\{x = 1\}}, will satisfy node coverage of this function.
      }

    \part[4] Is it possible to devise a set of test cases that achieve
      100\% \textbf{edge coverage} for \mintinline{java}{process}?  %
      If your answer is yes, provide a minimal set of test cases that
      do so. %
      If your answer is no, (a) provide a simple change to
      \mintinline{java}{process} so that it becomes possible to
      satisfy 100\% \textbf{edge coverage} and (b) provide a simple
      and minimal test case set that satisfies \textbf{edge coverage}
      of this function.

  \answerspace{14}
  \answer{\\
        % rubric:
        % Case 1 - Yes:
        %       - explain: 2
        %       - test case: 1
        % Case 2 - No:
        %       - explain: 1
        %       - change: 1
        %       - test case: 1
    No. Because within the while loop, \texttt{x} will never be less
    than zero due to \texttt{unsigned int} type.
    Thus, the node, \texttt{line 7}, will never be covered. \\
    Example change: To satisfy node coverage with a simple change,
    changing \texttt{unsigned int x} to \texttt{int x}, and adding
    \texttt{x -= 1;} after line $5$ in the while loop.  In this case,
    the test case, \texttt{\{x = 1\}}, will satisfy node coverage of
    this function.  }
  \end{parts}

  \newpage

\question[3] Is it possible for a set of test cases to achieve
  \textbf{complete path coverage} for a program? %
  If your answer is yes, provide an example of such a program and set
  of test cases. %
  If your answer is no, explain why not.


\question Consider the following faulty Java program:
  \begin{minted}{java}
// Sums all elements in a list
public static int sum(List<Integer> numbers) {
  int total = 0;
  for (int i = 0; i < numbers.size() - 1; i++) {
    total += numbers[i];
  }
  return total;
}
\end{minted}

\begin{parts}
\part[3] State the location(s) of the program fault and explain how
  you would fix the program. Consider only the smallest and simplest
  fix.
  \answerspace{10}
  \answer{\\
        % rubric:
        % location: 1
        % explain: 1
        % fix: 1
        %% Comment says the null pointer case is handled in exception.
        %One of the faults occurs in line 8, where the function does not check the a is null or not before the dereference. \\
        %Fix this fault by adding null checking before the for loop. \\
        Fault occurs in line 8, where the potentail out-of-bound access via the wrong condition in the for loop.
        Fix this fault by modifying \texttt{"i <= a.length"} to \texttt{"i < a.length"}. \\
    }

  \part[3] If possible, provide an example of a test input that does not cause the
    fault to execute. If this is not possible, explain why.
    \answerspace{10}
    \answer{\\
      % rubric:
        % Case 1 - have an example:
        %   - have: 1
        %   - input: 2
        % Case 2 - don't have an example:
        %   - don't have: 1
        %   - explain: 2
        It is impossible to avoid the fault in line 8,
        since there is no selection, jump, or iteration statements to let the program not to execute line 8.
    }

  \part[3] If possible, provide an example of a test input that causes the fault to execute
	but no failure to occur. If this is not possible, explain why.
	\answerspace{10}
    \answer{\\
        % rubric:
        % Case 1 - have an example:
        %   - have: 1
        %   - input satisfy requirement: 2
        % Case 2 - don't have an example:
        %   - don't have: 1
        %   - explain: 2
		% NOTE: Stating just the out-of-bounds or extra iteration is enough for full points on the justification.
        %If \texttt{a = null}, the program will throw a \texttt{NullPointerException}, which is expected behavior. This will execute the fault but no error to occur.
        %There is no possible way to have a test input that causes the fault to execute but no error to occur.
        %The program will throw an out-of-bounds access exception or the program will always execute an extra iteration if the language runtime does not detect it. (Either justification is acceptable).
    }

\end{parts}

\newpage

\question Consider the following faulty Java program:
  \begin{minted}{java}
    public static void main(String[] args) {
        Scanner scanner = new Scanner(System.in);

        int secretNumber = 42;
        int guess;
        int attempts = 0;

        System.out.println("Try to guess the secret number between 1 and 100.");

        while (true) {
            System.out.print("Enter your guess: ");
            guess = scanner.nextInt();
            attempts++;

            if (guess == secretNumber) {
                System.out.println("Correct! You guessed it in " + attempts + " attempts.");
                break;
            } else if (guess < secretNumber) {
                System.out.println("Too low!");
            } else {
                System.out.println("Too high!");
            }
            if (attempts == 5) {
                int hint = secretNumber \% 10
                System.out.println("Hint: The last digit of the secret number is " + hint);
            }
        }
        scanner.close();
    }
}
\end{minted}

\begin{parts}
\part[3] List all the DU-pairs for the variable
  \mintinline{java}{secretNumber}. Please use the format (Line\#X, Line\#Y).
  \answerspace{7}
  \answer{ }

\part[3] List all the DU-pairs for the variable
  \mintinline{java}{guess}. Please use the format (Line\#X, Line\#Y).
  \answerspace{7}
  \answer{ }

\part[3] List all the DU-pairs for the variable
  \mintinline{java}{attempts}. Please use the format (Line\#X, Line\#Y).
  \answerspace{7}
  \answer{ }
\end{parts}

\newpage

\question Consider the following pair of Java threads:
  \begin{minted}{java}
    Thread t1 = new Thread(() -> {
            A.lock();
            sharedCounter++;
            A.unlock();
            B.lock();
            sharedCounter--;
            B.unlock();
        });

        Thread t2 = new Thread(() -> {
            B.lock();
            sharedCounter += 2;
            B.unlock();
            A.lock();
            sharedCounter += 3;
            A.unlock();
        });
      \end{minted}
And suppose these threads execute according to the following schedule:
  \begin{enumerate}
  \item \textbf{T1:} L1-L3
  \item \textbf{T2:} L1-L6
  \item \textbf{T1:} L4-L6
  \end{enumerate}
\textbf{Answer format:} When writing candidate locksets or held locks, always use curly braces
(\texttt{\{\}}).
For example: \texttt{\{\}} (empty set), \texttt{\{A\}}, \texttt{\{A, B\}}. Do not write \texttt{None}, \texttt{empty} or leave the text field empty.\\

\begin{parts}

\part[2] Candidate lockset for \mintinline{java}{sharedCounter} before step 2: (Use the \textbf{answer format} specified above)
\answerspace{10}
\answer{
	\{\}
}
\part[2] Candidate lockset for \mintinline{java}{sharedCounter} between steps 2+3: (Use the \textbf{answer format} specified above)
\answerspace{10}
\answer{
	\{\}
}
\part[2] Candidate lockset for \mintinline{java}{sharedCounter} after step 3: (Use the \textbf{answer format} specified above)
\answerspace{10}
\answer{
	\{A\}
}
\part[2] Does the lockset data race algorithm flag this execution as
  containing a potential data race?
\answerspace{10}
\answer{

}
\part[4] Is this an actual data race and, if so, is it a benign data
  race? Explain your answer.
\answerspace{10}
\answer{
	Yes. T2 writes to \texttt{x} without holding \texttt{A} while T1 writes while holding \texttt{A}.
}
\end{parts}

\newpage

\question Consider the following Java program:
  \begin{minted}{java}
public enum StudentCategory {
  grad, honors, undergrad
};

double calculateCurve(StudentCategory cat, double origScore) {
  case grad:
    return origScore;
  case honors:
    return (origScore / 2) + 50;
  default:
    return (origScore / 2) + 53;
}
\end{minted}
and the following set of blocks:
\begin{tabular}{c|c|c}
  \textbf{block} & \textbf{status} & \textbf{grade} \\
  b1 & graduate & F \\
  b2 & graduate & below C \\
  b3 & honors & below D \\
  b4 & undergrad & any \\
\end{tabular}

\begin{parts}
\part[5] Is this a valid partitioning of the input space of
  \mintinline{java}{calculateCurve}? Explain your answer.

\part[5] If your answer was no, how would you modify and/or augment
  this set of blocks to turn it into a valid partitioning?
\end{parts}

\if 0 %%%%%%%%%%%%%%%%%%%%%%%%%%%%%%%%%%%%%%%%%
	\question[10]
	\answerspace{30}
	\answer{}

	\question
	\begin{parts}
		\part[6]
		\answerspace{20}
		\answer{}

		\part[2]
		\answerspace{10}
		\answer{}

	\end{parts}
\fi %%%%%%%%%%%%%%%%%%%%%%%%%%%%%%%%%%%%%%%%%

	\newpage
	\begin{center}
		[EMPTY PAGE 1]
	\end{center}

	\newpage
	\begin{center}
		[EMPTY PAGE 2]
	\end{center}

\end{questions}





\end{document}
